\documentclass[10pt,a4paper]{amsart}
\usepackage[margin=19mm]{geometry}
\usepackage{amsmath,amssymb,mathtools}
\usepackage[scr=rsfs,cal=euler]{mathalfa}
\usepackage{xfrac}
\usepackage[unicode,hidelinks,bookmarks=false]{hyperref}
\newcommand{\ie}{i.e.\ }
\newcommand{\cf}{cf.\ }
\newcommand{\resp}{respectively\ }
\newcommand{\Subsection}{\subsection}
\providecommand{\nobreakdash}{\nobreak\hbox{-}}
\setlength{\emergencystretch}{2em}
\multlinegap0pt
\usepackage{bm}
\usepackage{bbm}
\usepackage{dsfont}
\usepackage{xspace}
\usepackage{cancel}
\usepackage{mathtools}
\usepackage[shortlabels]{enumitem}
\usepackage{amsthm}
\makeatletter
\@ifundefined{equation}{\newtheorem{equation}{Equation}}{}
\@ifundefined{Prop}{\newtheorem{Prop}[equation]{Proposition}}{}
\makeatother
\newcommand{\NN}{\mathbb{N}}
\newcommand{\ZZ}{\mathbb{Z}}
\title[Topologizing infinite quivers and their mutations]{Topologizing infinite quivers and their mutations}
\author{Benjamin Grant}
\date{Source: arXiv:2604.16660v1 (CC BY 4.0)}
\begin{document}
\maketitle

\noindent\emph{Source and licence.} Topologizing infinite quivers and their mutations. Version arXiv:2604.16660v1 (CC BY 4.0), distributed under the Creative Commons Attribution 4.0 International licence, as recorded in the arXiv licence metadata for this exact version and shown on the abstract page https://arxiv.org/abs/2604.16660v1. Full rights notice: licenses/arxiv-2604.16660-CC-BY-4.0.txt.\par\smallskip

\noindent\textit{Editorial scope.} Counted unit: the Prop whose source label is \texttt{linking-prop} in the article (source0.tex lines 1238--1246, source label \texttt{linking-prop}), delivered together with the article's own complete original proof of it (source0.tex lines 1247--1306, one \texttt{proof} environment of 60 lines). No mathematics is written, repaired or re-derived: statement and proof are the article's own text line for line. Required context retained uncounted: none. Every definition, display and label of those lines is retained unchanged. Omitted from this package and not redistributed: 1--1237, 1307--1754. Nothing inside a retained excerpt is omitted or altered: every retained body is delivered exactly as the article's lines are written, so no omission receipt of any kind accompanies this package. Citations: the retained excerpts carry no citation command at all, so no bibliography entry of the article is transcribed and no reference is redistributed. Reference adaptation: none was needed and none was made; every reference inside the delivered unit resolves inside it, so no unresolved reference or citation is produced. Printed slips kept literal: no printed word, symbol, hypothesis or number of the counted statement or of its proof was altered. Layout and class: the article's own class is \texttt{amsart} with packages including latexsym, mathrsfs, color, tikz, multirow, bbm, mathtools, amsmath, amssymb, enumitem, quiver, tikz-cd, circuitikz, dotlessi; the wrapper is the lane's pinned \texttt{amsart} editorial wrapper at 10pt on A4, which restates the theorem-like environments the retained text needs and adds the article's own macro definitions that the retained text needs (\texttt{\textbackslash NN}, \texttt{\textbackslash ZZ}), with extra wrapper packages (bm, bbm, dsfont, xspace, cancel, mathtools, enumitem). The delivered numbering is the unit's own. Assets: no retained span contains a figure environment, a graphics-inclusion command, a PDF-page inclusion, a table or a TikZ picture; no image file is copied into this package, the package declares no asset dependency and no image file is redistributed with it. Licence: version arXiv:2604.16660v1 of this article is distributed under the Creative Commons Attribution 4.0 International licence, as recorded in the arXiv licence metadata for this exact version and shown on the abstract page https://arxiv.org/abs/2604.16660v1. A full rights notice is delivered at licenses/arxiv-2604.16660-CC-BY-4.0.txt. Characters: every retained excerpt is pure ASCII, so the delivered engine, XeLaTeX with its default Latin Modern text font, reports no missing character.
\begin{Prop}\label{linking-prop}
    Let $\mathbf i\in\NN^\NN$ be such that each $i\in\NN$ appears only finitely many times in $\mathbf i$ and so that for all $k\geq1$, it holds that $i_k\neq i_{k+1}$ (equivalently, $\mathbf i=R(\mathbf i)=R^\omega(\mathbf i)$ is reduced). Then the following four properties hold of $\mathcal L_\mathbf i$:
    \begin{enumerate}[label=(\roman*)]
        \item For any nonempty finite set $\varnothing\neq S\subseteq\NN$, there are only finitely many finite sets $S'\supseteq S$ such that $S'\in\mathcal L_\mathbf i$ \emph{and} $S'$ is minimal (with respect to inclusion) in $\mathcal L_\mathbf i$.
        \item There exist infinitely many finite sets $S\in\mathcal L_\mathbf i$ which are minimal in $\mathcal L_\mathbf i$.
        \item For every $S\in\mathcal L_\mathbf i$, there is some finite $S'\subseteq S$ belonging to $\mathcal L_\mathbf i$ which is minimal in $\mathcal L_\mathbf i$.
        \item All minimal $S\in\mathcal L_\mathbf i$ are finite.
    \end{enumerate}
\end{Prop}

\begin{proof}
    We work toward (i) first. We claim that it suffices to show that this condition holds for $S=\{i\}\subseteq\NN$ a singleton; suppose under the assumption that (i) holds for singletons $S$ that there is some finite $T\subseteq\NN$ contained in infinitely many minimal finite sets $S'\in\mathcal L_\mathbf i$. Then for any $i\in T$, the singleton $\{i\}$ is contained in infinitely many such $S'$ as well, a contradiction. Therefore, it suffices to show that (i) holds for singletons.

    To proceed, we define the \emph{convex hull of $i\in\NN$ with respect to $\mathbf i$} to be the set $\text{conv}_\mathbf i(i)\subseteq\NN$ given by
    $$\text{conv}_\mathbf i(i)=\left\{j\in\NN\mid\exists k_1\leq k\leq k_2:i=i_{k_1}=i_{k_2}\text{ and }j=i_k\right\}.$$
    Note that $\text{conv}_\mathbf i(i)=\varnothing$ if $i$ does not appear in $\mathbf i$ and $i\in\text{conv}_\mathbf i(i)$ otherwise. Observe further that $\text{conv}_\mathbf i(i)$ is finite by the assumption that no $i\in\NN$ appears in $\mathbf i$ infinitely many times. We claim two important things: that $\text{conv}_\mathbf i(i)\in\mathcal L_\mathbf i$ for all $i$ appearing in $\mathbf i$ and, moreover, that all finite minimal sets $S'\in\mathcal L_\mathbf i$ containing $i$ are subsets of $\text{conv}_\mathbf i(i)$. If we can prove these, then we have shown part (i) of the proposition: there are only finitely many $S'\subseteq\text{conv}_\mathbf i(i)$, and therefore only finitely many finite \emph{minimal} $S'\in\mathcal L_\mathbf i$ containing $i$.

    To argue toward either of these two claims, we will appeal to the group $G=\langle g_i\mid g_i^2=\varepsilon,i\in\NN\rangle$ and certain endomorphisms of $G$. Elements of $G$ can be naturally identified with finite sequences $(j_1,j_2,\dots,j_\ell)\in\NN^{<\NN}$ with $\ell\geq0$ such that $j_k\neq j_{k+1}$ for all $1\leq k<\ell$: simply take such a $\mathbf j=(j_1,j_2,\dots,j_\ell)$ to the group element $g_\mathbf j=g_{j_1}g_{j_2}\dots g_{j_\ell}\in G$. Note that $g_\mathbf j=g_{R^\omega(\mathbf j)}$; the group $G$ ``automatically'' performs reduction of sequences. For each $i\in\NN$, define an endomorphism $\varphi_i:G\to G$ on the generators of $G$ by
    $$\varphi_i(g_j)=\begin{cases}
        \varepsilon & \text{if }j=i\\
        g_j & \text{if }j\neq i.
    \end{cases}$$
    We note the following useful property: for any $g\in G$ and $i\in\NN$, it holds that $\varphi_i(g)=g$ if and only if $g$ is in the image of $\varphi_i$. Certainly the ``only if'' direction holds, so suppose that $g$ is in the image of $\varphi_i$, say $g=\varphi_i(h)$ for $h\in G$. Since $\varphi_i$ is idempotent, we have that $\varphi_i(g)=\varphi_i^2(h)=\varphi_i(h)=g$.
    
    Another useful property is that for $\mathbf j=(j_1,\dots,j_\ell)\in\NN^{<\NN}$ with $j_k\neq j_{k+1}$ for all $1\leq k<\ell$ as above and for any $i\in\NN$, it holds that $i$ appears in $\mathbf j$ if and only if $g_\mathbf j$ is not in the image of $\varphi_i$. We argue this using the \emph{length function} $L:G\to\ZZ_{\geq0}$ given by $L(g_{(j_1,\dots,j_\ell)}):=\ell$. Note that two facts hold of the length function $L$: that for all $g,h\in G$, we have $L(gh)\leq L(g)+L(h)$ and that for all $g\in G$, $i\in\NN$, we have $L(\varphi_i(g))\leq L(g)$. Now suppose first that $\mathbf j$ contains an instance of $i$, and let $\mathbf j':=\mathbf j_{\NN\setminus\{i\}}$ be the subsequence of $\mathbf j$ obtained by removing all instances of $i$. This yields
    $$L(\varphi_i(g_\mathbf j))=L(g_{\mathbf j'})=|R^\omega(\mathbf j')|<|\mathbf j|=L(g_\mathbf j),$$
    so $\varphi_i(g_\mathbf j)\neq g_\mathbf j$. By our previous useful property, this implies that $g_\mathbf j$ is not in the image of $\varphi_i$. For the converse, suppose that $\mathbf j$ does not contain any instances of $i$. Then, by the definition of $\varphi_i$, we see that $\varphi_i(g_\mathbf j)=g_\mathbf j$, so $g_\mathbf j$ is in the image of $\varphi_i$. Altogether, this gives us a group-theoretic criterion for detecting when an element $i\in\NN$ appears in a reduced finite sequence $\mathbf j$.

    We now establish that $\text{conv}_\mathbf i(i)$ is linked with respect to $\mathbf i$ for $i$ appearing in $\mathbf i$. Observe that we may write $\gamma:=g_{\mathbf i_{\text{conv}_\mathbf i(i)}}=\tau'\sigma\tau$ for $\tau',\sigma,\tau\in G$ such that:
    \begin{enumerate}
        \item $\sigma\neq\varepsilon$, i.e. $\sigma$ comes from a substring of $\mathbf i_{\text{conv}_\mathbf i(i)}$ with nonempty reduction.
        \item $\varphi_i(\sigma)\neq\sigma$, i.e. $g_i$ appears in the reduced expression for $\sigma$.
        \item $\varphi_i(\tau)=\tau$, $\varphi_i(\tau')=\tau'$, i.e., $g_i$ does not appear in the reduced expressions for $\tau$ or $\tau'$.
    \end{enumerate}
    Indeed, we may take $\sigma=g_\mathbf j$ for the consecutive substring $\mathbf j=(i_{e_1},i_{e_1+1},\dots,i_{e_2-1},i_{e_2})$ of $\mathbf i$, where $i_{e_1}$ is the first appearance of $i$ in $\mathbf i$ and $i_{e_2}$ is the last. We then take $\tau'$ to capture all instances of any $j\in\text{conv}_\mathbf i(i)$ appearing at indices less than $e_1$ and we take $\tau$ to capture all instances of any $j\in\text{conv}_\mathbf i(i)$ appearing at indices after $e_2$.
    
    We claim that $\gamma\neq\varepsilon$, or equivalently that $R^\omega(\mathbf i_{\text{conv}_\mathbf i(i)})\neq\varnothing$. Suppose for the sake of contradiction that $\gamma=\varepsilon$. Rewriting $\gamma=\tau'\sigma\tau$, we then see that $\sigma=(\tau\tau')^{-1}$. Applying $\varphi_i$ to both sides yields
    $$\varphi_i(\sigma)=\varphi_i((\tau\tau')^{-1})=(\tau\tau')^{-1}=\sigma,$$
    a contradiction with $\varphi_i(\sigma)\neq\sigma$.

    Toward the latter of our two important claims, we conduct a similar argument to show that all finite minimal sets $S'\in\mathcal L_\mathbf i$ are subsets of $\text{conv}_\mathbf i(i)$. Suppose for the sake of contradiction that $S'\in\mathcal L_\mathbf i$ is a finite minimal set containing $i$ which is not contained in $\text{conv}_\mathbf i(i)$; take some $j\in S'\setminus \text{conv}_\mathbf i(i)$. We again may write $\gamma:=g_{\mathbf i_{S'}}=\tau'\sigma\tau$ for $\tau',\sigma,\tau\in G$ such that:
    \begin{enumerate}
        \item $\sigma\neq\varepsilon$, i.e. $\sigma$ comes from a substring of $\mathbf i_{\text{conv}_\mathbf i(i)}$ with nonempty reduction.
        \item $\varphi_i(\sigma)\neq\sigma$, i.e. $g_i$ appears in the reduced expression for $\sigma$.
        \item $\varphi_i(\tau)=\tau$, $\varphi_i(\tau')=\tau'$, i.e., $g_i$ does not appear in the reduced expressions for $\tau$ or $\tau'$.
        \item $\varphi_j(\sigma)=\sigma$, i.e., $g_j$ does not appear in the reduced expression for $\sigma$.
        \item Either $\varphi_j(\tau')\neq\tau'$ or $\varphi_j(\tau)\neq\tau$ (or both). That is, there is an instance of $g_j$ in at least one of the reduced expressions for $\tau$ and $\tau'$. We may assume without loss of generality that $\varphi_j(\tau')\neq\tau'$.
        \item $\gamma\neq\varepsilon$, i.e., $R^\omega(\mathbf i_{S'})\neq\varnothing$.
    \end{enumerate}
    Indeed, condition (4) here is capturing the fact that $j\not\in\text{conv}_\mathbf i(i)$, while condition (5) is requiring our use of at least one instance of $j$ in $\mathbf i_{S'}$. We also may reformulate the minimality condition on $S'\in\mathcal L_\mathbf i$ in the following group-theoretic fashion: it is equivalent to say that $\gamma\in\bigcap_{s\in S'}\ker\varphi_s$. Indeed, taking away all instances of $s\in S'$ in $\mathbf i_{S'}$ will yield a sequence $\mathbf i_{S'\setminus\{s\}}$ which reduces to $\varnothing$ by the minimality of $S'$, so $\varphi_s(\gamma)=\varepsilon$ for all $s\in S'$.

    We now obtain a contradiction in a similar manner to before. Take $\varphi_j$ on both sides of $\gamma=\tau'\sigma\tau$ to see that
    \begin{align*}
        \varepsilon&=\varphi_j(\tau'\sigma\tau)\\
        &=\varphi_j(\tau')\varphi_j(\sigma)\varphi_j(\tau)\\
        &=\varphi_j(\tau')\sigma\varphi_j(\tau),\\
    \end{align*}
    yielding $\sigma=\varphi_j(\tau\tau')^{-1}$. Then, applying $\varphi_i$ to both sides of this equation (and noting that $\varphi_i\circ\varphi_j=\varphi_j\circ\varphi_i$) yields
    \begin{align*}
        \varphi_i(\sigma)&=\varphi_i(\varphi_j(\tau\tau'))^{-1}\\
        &=\varphi_j(\varphi_i(\tau\tau'))^{-1}\\
        &=\varphi_j(\tau\tau')^{-1}\\
        &=\sigma,
    \end{align*}
    a contradiction to condition (2) above. Thus, we cannot have a finite minimal $S'\in\mathcal L_\mathbf i$ containing $i$ which is not contained in $\text{conv}_\mathbf i(i)$. This concludes the proof of part (i).

    Part (ii) follows from part (i) and the fact that there are infinitely many $i\in\NN$ appearing in $\mathbf i$. Indeed, for any $i\in\NN$ appearing in $\mathbf i$, by the argument for part (i) provided above, there exists at least one finite minimal set $S'\in\mathcal L_\mathbf i$ containing $i$, and there exist only finitely many such $S'$. Let $m\in\NN$ be the maximum element (in the natural order on $\NN$) of all such $S'$. Repeat this argument for the least $i'>m$ appearing in $\mathbf i$ to obtain more finite minimal sets; repeat this argument ad infinitum to obtain infinitely many.

    Finally, for parts (iii) and (iv), we first observe that part (iv) is an immediate consequence of part (iii), so this is all that is left to show. Let $S\in\mathcal L_\mathbf i$. We want to show there is some finite minimal $S'\subseteq S$ which is itself linked with respect to $\mathbf i$. Note that it suffices to show that there is some finite $S'\subseteq S$ which is linked. Assume without loss of generality that $S$ is infinite. If $R^\omega(\mathbf i_S)\neq\varnothing$ is finite, simply let $S'$ be the finite set of $j\in\NN$ appearing in $R^\omega(\mathbf i_S)$. Then $S'$ is linked in $\mathbf i$ precisely because $R^\omega(\mathbf i_{S'})=R^\omega(\mathbf i_S)$. Hence, we may suppose that $R^\omega(\mathbf i_S)$ is infinite; replace $\mathbf i_S$ with $R^\omega(\mathbf i_S)$ without loss of generality. Now $\mathbf i_S$ is an infinite reduced sequence with no $i\in\NN$ appearing infinitely many times. In particular, we may apply our argument for part (i) to $\mathbf i_S$ to obtain for any $i\in S$ a finite set $\text{conv}_{\mathbf i_S}(i)\subseteq S$ containing $i$ which is linked with respect to $\mathbf i_S$. Then it also follows that $\text{conv}_{\mathbf i_S}(i)$ is linked with respect to the original sequence $\mathbf i$, so it holds that $S':=\text{conv}_{\mathbf i_S}(i)\in\mathcal L_\mathbf i$ meets our conditions. This completes the proof.
\end{proof}

\end{document}
